\scanpage{31}
% archive sheet number: 30
\mspage{- 47 -}\hfill n$^{\mathrm{o}}$ 272

$n=\handwritten{[}L:K\handwritten{]}$. \quad b) \uline{L'intersection} $B$ \uline{des} $A_i$ \uline{est la cl\handwritten{\^}{o}ture intégrale de} $A$ dans $L$.

c) \uline{Si l'on pose} $\underline{p}_i=B\handwritten{\cap}\underline{m}(A_i)$, \uline{les idéaux} $\underline{p}_i$ \uline{sont des idéaux maximaux de} $B$, \uline{distincts} : \uline{tout idéal maximal de} $B$ \uline{est égal à l'un des} $\underline{p}_i$ \uline{et l'on a} \unsure{.|} $A_i=B_{\underline{p}_i}$.

Nous noterons $v_i$ (resp.$v$) des valuations associées aux $A_i$ (resp. à $A$), et $\Gamma_i$ (\unsure{r\.esp.|rèsp.|resp.}$\Gamma$) les groupes des ordres correspondants. Le groupe $\Gamma$ s'identifie à un sous-groupe de chacun des groupes $\Gamma_i$.

Montrons d'abord que l'on ne peut pas avoir $A_i\handwritten{\supset}A_j$ si $i\handwritten{\not}=j$. En effet, d'après ce que l'on a vu au \handwritten{\S}4, le groupe $\Gamma_i$ s'identifierait à un groupe quotient $\Gamma_j\handwritten{/}H$, où $H$ est un sous-groupe isolé de $\Gamma_j$ ; puisque l'application canonique $\Gamma\handwritten{\longrightarrow}\Gamma_j\handwritten{\longrightarrow}\Gamma_j\handwritten{/}H=\Gamma_i$ est une injection, on a $\handwritten{\Gamma\cap}H=\handwritten{\{}0\handwritten{\}}$ dans $\Gamma_j$, ce qui\unsure{.|...|} montre que $H$ est isomorphe à un sous-groupe de $\Gamma_j\handwritten{/}\Gamma$ ; comme ce dernier groupe est fini (prop.1), on a $H=0$, et les valuations $v_i$ et $v_j$ sont équivalentes, d'où $A_i=A_j$ contrairement à l'hypothèse $i\handwritten{\not}=j$.

Montrons maintenant que le nombre des $A_i$ est \unsure{1|,|'} $\handwritten{\leqslant}n$. Supposons donc qu'il existe $s$ anneaux $A_i$ distincts, avec $s\handwritten{>}n$. D'après le cor. 2 à la \unsure{.|}prop.2 du \handwritten{\S}6, (qui est applicable, en vertu de ce que l'on vient de voir), il existe des éléments $x_i\handwritten{\in}L$, $1\handwritten{\leqslant}i\handwritten{\leqslant}s$, tels que $v_i(x_i)=0$ et $v_j(x_i)\handwritten{>}0$ pour $i\handwritten{\not}=j$. Montrons que les $x_i$ sont linéairement indépendants sur $K$, ce qui fournira la contradiction cherchée. Supposons que l'on ait une relation linéaire non triviale :
\[
\handwritten{\sum} a_i x_i=0\quad,\ \text{avec }a_i\handwritten{\in}K.
\]

Soit $a_j$ celui des $\unsure{a'_i|a_i}$ dont la valuation est minimum ; quitte à multiplier les $a_i$ par $a_j^{-1}$ on peut supposer que $a_j=1$. On a alors une relation :

